\documentclass[11pt]{article}
\usepackage{preamble}
\setlength{\parskip}{4pt}

\begin{document}

\daybanner{AP Statistics \textperiodcentered\ Unit 3 \textperiodcentered\ Topic 3.5 \textperiodcentered\ B: 2026-12-03 \textperiodcentered\ E: 2027-01-08 \textperiodcentered\ TEACHER KEY}{Which Test Answers the Question Asked?}

\sectionbanner{CED TETHER}

{\small
\begin{itemize}[leftmargin=1.5em,itemsep=2pt,topsep=2pt]
  \item \textbf{Enduring Understanding VAR-6} --- The normal distribution may be used to model variation.
  \item \textbf{LO VAR-6.D} --- Identify the null and alternative hypotheses for a population proportion. [Skill 1.F]
  \item \textbf{VAR-6.D.1} --- The null hypothesis is the situation that is assumed to be correct unless evidence suggests otherwise, and the alternative hypothesis is the situation for which evidence is being collected.
  \item \textbf{VAR-6.D.2} --- For hypotheses about parameters, the null hypothesis contains an equality reference (=, $\geq$, or $\leq$), while the alternative hypothesis contains a strict inequality (<, >, or $\neq$). The type of inequality in the alternative hypothesis is based on the question of interest. Alternative hypotheses with < or > are called one-sided, and alternative hypotheses with $\neq$ are called two-sided. Although the null hypothesis for a one-sided test may include an inequality symbol, it is still tested at the boundary of equality.
  \item \textbf{VAR-6.D.3} --- The null hypothesis for a population proportion is: $H_0$: p = $p_0$, where $p_0$ is the null hypothesized value for the population proportion.
  \item \textbf{VAR-6.D.4} --- A one-sided alternative hypothesis for a proportion is either $H_a$: p < $p_0$ or $H_a$: p > $p_0$. A two-sided alternative hypothesis is $H_a$: p $\neq$ $p_0$.
  \item \textbf{VAR-6.D.5} --- For a one-sample z-test for a population proportion, the null hypothesis specifies a value for the population proportion, usually one indicating no difference or effect.
  \item \textbf{LO VAR-6.E} --- Identify an appropriate testing method for a population proportion. [Skill 1.E]
  \item \textbf{VAR-6.E.1} --- For a single categorical variable, the appropriate testing method for a population proportion is a one-sample z-test for a population proportion.
  \item \textbf{LO VAR-6.F} --- Verify the conditions for making statistical inferences when testing a population proportion. [Skill 4.C]
  \item \textbf{VAR-6.F.1} --- In order to make statistical inferences when testing a population proportion, we must check for independence and that the sampling distribution is approximately normal:
  \item a. To check for independence:
  \item i. Data should be collected using a random sample or a randomized experiment.
  \item ii. When sampling without replacement, check that n $\leq$ 10\%N.
  \item b. To check that the sampling distribution of p-hat is approximately normal (shape):
  \item i. Assuming that $H_0$ is true (p = $p_0$), verify that both $np_0$ and n(1 - $p_0$) are at least 10 so that the sample size is large enough to support an assumption of normality.
\end{itemize}
}
\textbf{Essential question:} How do the research direction and null model determine a valid one-proportion test setup?\par
\textbf{DOK-3 skill:} 4.C \textperiodcentered\ \textbf{FRQ pattern:} {\small\texttt{directional-\allowbreak{}question-\allowbreak{}null-\allowbreak{}based-\allowbreak{}conditions-\allowbreak{}which-\allowbreak{}test-\allowbreak{}setup}}\par
\textbf{DOK rationale:} Part (c) adjudicates two plausible test setups by coordinating the research direction, parameter-versus-statistic distinction, procedure choice, and null-based condition checks, then explains what each mismatch would test instead.\par

\frameworkphaseheader{Do Now (first take) $\rightarrow$ Explore (video + follow-along) $\rightarrow$ Exit (finish + turn in)}{1 $\rightarrow$ 3}{45}{%
\begin{itemize}[leftmargin=1.5em,itemsep=2pt,topsep=2pt]
  \item Board slide up at the bell. Ask students to choose a setup before confirming it.
  \item Begin the video at minute 5; return at minute 33 to label each number's role.
  \item Collect at the door. Score part (c) only, E/P/I.
\end{itemize}
}{%
\begin{itemize}[leftmargin=1.5em,itemsep=2pt,topsep=2pt]
  \item Choose Maya or Theo and cite the direction or a condition.
  \item Complete the follow-along, finish (a)--(c), revise, and turn in the sheet.
\end{itemize}
}{%
\begin{itemize}[leftmargin=1.5em,itemsep=2pt,topsep=2pt]
  \item What do $p$ and $84/180$ each describe?
  \item Would a lower current sample proportion count as evidence that ridership increased?
  \item Does large counts model the null or the observed sample?
\end{itemize}
}{Solo teacher. Circulate ~5 pairs. Connect every choice to the question or null model.}

\begin{callout}{\IconWarn\ WATCH FOR}{calloutred}
\begin{itemize}[leftmargin=1.5em,itemsep=2pt,topsep=2pt]
  \item Choosing two-sided merely because samples vary both ways.
  \item Putting the sample proportion in a hypothesis.
  \item Using observed rather than null-expected counts.
\end{itemize}
\end{callout}

\sectionbanner{SCORING PART (c) --- E / P / I}

\textbf{Expected elements}
\begin{itemize}[leftmargin=1.5em,itemsep=2pt,topsep=2pt]
  \item \textbf{judgment-1} (required) --- Chooses Maya and ties the greater-than alternative to the pre-stated increase.
  \item \textbf{judgment-2} (required) --- Explains that 72 and 108 are expected under the null benchmark p=0.40, whereas 84 and 96 are observed counts.
  \item \textbf{judgment-3} (required) --- Explains that Theo's two-sided alternative would treat a decrease as evidence of change.
\end{itemize}

{\small\renewcommand{\arraystretch}{1.25}
\noindent\begin{tabular}{|p{0.5in}|p{5.9in}|}
\hline
\textbf{E} & Meets all three checks with correct contextual reasoning; equivalent wording is acceptable. \\ \hline
\textbf{P} & Shows at least one substantive correct check, but has an omission or error that prevents E. \\ \hline
\textbf{I} & Shows none of the three checks with substantive correct reasoning; a choice alone is insufficient. \\ \hline
\end{tabular}}

\clearpage
\focusbanner{TODAY'S PROBLEM --- annotated}

A district of 2{,}400 students had bus ridership $p_0=0.40$ before a new route. Before collecting data, the director asks whether the current proportion $p$ has \emph{increased}. An SRS of 180 current students finds 84 riders.\par\smallskip \textbf{Maya:} ``Use $H_0:p=0.40$ versus $H_a:p>0.40$ and a one-sample $z$-test. Check large counts with 72 and 108.''\par \textbf{Theo:} ``Use $H_a:p\neq0.40$ because samples can differ either way. Check large counts with the observed 84 and 96.''\par
\begin{center}
\textbf{Current SRS counts}\par\smallskip
\begin{tabular}{|l|c|c|c|}
\hline
\textbf{Group} & \textbf{Bus} & \textbf{Other} & \textbf{Total} \\ \hline
\textbf{Students} & 84 & 96 & 180 \\ \hline
\end{tabular}
\end{center}
\begin{firsttakebox}{FIRST TAKE (5 min)}
Does the director want to detect any change, or only more bus riders? Use a phrase from the study.\par
\answer{Not graded. Separate the question asked from the sample observed.}
\end{firsttakebox}

\textit{Rules callout ``SET UP THE QUESTION BEFORE USING THE SAMPLE (keep on the page)'' is printed on the student sheet.}\par\medskip

\dokbadge{1}~\textbf{(a)} Define the population parameter $p$ in context and identify the null value $p_0$.\par
\answer{Here $p$ is the current proportion of all 2{,}400 district students who use the school bus. The null value is the pre-route benchmark $p_0=0.40$.}

\dokbadge{2}~\textbf{(b)} Compute $\hat p$, name the testing procedure, and verify the random, 10 percent, and large-counts conditions using the appropriate values.\par
\answer{$\hat p=84/180=7/15\approx0.4667$. Use a one-sample $z$-test for $p$. The SRS is random; $180\le0.10(2{,}400)=240$; and under $H_0$, $180(0.40)=72$ and $180(0.60)=108$ are at least 10. All conditions hold.}

\dokbadge{3}~\textbf{(c)} In 3--4 sentences, choose a setup using the director's word increased. Explain why the count check uses 72 and 108 rather than 84 and 96. State what different question Theo's alternative would answer.\par
\answer{Maya is warranted. The pre-data word \emph{increased} requires $H_a:p>0.40$; $p_0=0.40$ is the benchmark, while $\hat p=84/180$ is evidence, not a hypothesis value. Large counts models $H_0$, so it uses expected counts 72 and 108. Theo's two-sided test asks about any change and would count a decrease as evidence. His observed 84 and 96 happen to exceed 10, but using them could make validity change with the sample and mislead readers about the null sampling distribution.}

\begin{summaryexitbox}{EXIT}
One thing the video changed about my first take: \hrulefill
\end{summaryexitbox}

\end{document}
